\label{sec:extremos}

\begin{definition}
  Sea $f:U\subseteq\Rn{n}\to\R$ de clase $\mathcal{C}^1$. Se define la \emph{matriz hessiana de $f$ en $\mathbf{x}_0$}, y se nota por $\boldsymbol{H}_f(\mathbf{x}_0)$, a 
  \[
    \boldsymbol{H}_f(\mathbf{x}_0)=
    \begin{bmatrix}
      \grad{\dif{f}{x_1}(\mathbf{x}_0)}\\[.4cm]
      \grad{\dif{f}{x_2}(\mathbf{x}_0)}\\[.4cm]
      \vdotswithin{f}\\[.4cm]
      \grad{\dif{f}{x_n}(\mathbf{x}_0)}
    \end{bmatrix}
    =
    \begin{bmatrix}
      \dfrac{\partial^2 f}{\partial x_1^2} & \dfrac{\partial^2 f}{\partial x_1 \partial x_2} &  \cdots & \dfrac{\partial^2 f}{\partial x_1 \partial x_n} \\[.4cm]
      \dfrac{\partial^2 f}{\partial x_2 \partial x_1} & \dfrac{\partial^2 f}{\partial x_2^2} & \cdots & \dfrac{\partial^2 f}{\partial x_2 \partial x_n} \\[.4cm]
      \vdots & \vdots & \ddots & \vdots \\[.4cm]
      \dfrac{\partial^2 f}{\partial x_n \partial x_1} & \dfrac{\partial^2 f}{\partial x_n \partial x_2} & \cdots & \dfrac{\partial^2 f}{\partial x_n^2}.
    \end{bmatrix}
\]
\end{definition}

\begin{obs}
Notar que $\boldsymbol{H}_f(\mathbf{x}_0)$ es una matriz cuadrada y sim\'etrica $\in \R^{n\times n}$, ya que $f\in\mathcal{C}^2$. \textcolor{red}{esta escrito el teorema de schwartz? donde lo pongo?}
\end{obs}

\begin{theorem} \textbf{Polinimio de Taylor de primer orden.}
  Sea $f:U\subseteq\Rn{n}\to\R$ de clase $\mathcal{C}^1$. Se define el \emph{polinomio de Taylor de primer orden centrado en $\mathbf{x}_0$ de $f$}, y se nota por $P_1[f,\mathbf{x}_0](\mathbf{x})$, a
  \[
    P_1[f,\mathbf{x}_0](\mathbf{x})=f(\mathbf{x}_0)+\grad{f}(\mathbf{x}_0)\cdot(\mathbf{x-x}_0).
  \]
  El \emph{resto del polinomio de Taylor de primer orden centrado en $\mathbf{x}_0$ de $f$}, y se nota por $R_1[f,\mathbf{x}_0](\mathbf{x})$, a
  \[
    R_1[f,\mathbf{x}_0](\mathbf{x})=f(\mathbf{x})-P_1[f,\mathbf{x}_0](\mathbf{x}).
  \]
  Entonces se cumple que
  \[
    \lim_{\mathbf{x-x}_0}\dfrac{R_1[f,\mathbf{x}_0](\mathbf{x})}{\|\mathbf{x-x}_0\|}=0.
  \]
  Si, adem\'as, $f\in\mathcal{C}^2\implies\exists\:\mathbf{c}\in\Rn{n}$, en el segmento que une $\mathbf{x}$ con $\mathbf{x}_0$, tal que \footnote{El super\'indice T denota la transpuesta de una matriz. Notar que se toma al vector $\mathbf{x}\in\R^{1\times n}$.}
  \[
    R_1[f,\mathbf{x}_0](\mathbf{x})=\frac{1}{2!}(\mathbf{x-x}_0)\boldsymbol{H}_f(\mathbf{c})(\mathbf{x-x}_0)^\text{T}.
  \]
  \textcolor{red}{definir segmento entre puntos?}
\end{theorem}

\begin{theorem} \textbf{Polinimio de Taylor de segundo orden.}
  Sea $f:U\subseteq\Rn{n}\to\R$ de clase $\mathcal{C}^2$. Se define el \emph{polinomio de Taylor de segundo orden centrado en $\mathbf{x}_0$ de $f$}, y se nota por $P_2[f,\mathbf{x}_0](\mathbf{x})$, a
  \[
    P_2[f,\mathbf{x}_0](\mathbf{x})=f(\mathbf{x}_0)+\grad{f}(\mathbf{x}_0)\cdot(\mathbf{x-x}_0)+\frac{1}{2!}(\mathbf{x-x}_0)\boldsymbol{H}_f(\mathbf{x}_0)(\mathbf{x-x}_0)^\text{T}.
  \]
  El \emph{resto del polinomio de Taylor de segundo orden centrado en $\mathbf{x}_0$ de $f$}, y se nota por $R_2[f,\mathbf{x}_0](\mathbf{x})$, a
  \[
    R_2[f,\mathbf{x}_0](\mathbf{x})=f(\mathbf{x})-P_2[f,\mathbf{x}_0](\mathbf{x}).
  \]
  Entonces se cumple que
  \[
    \lim_{\mathbf{x-x}_0}\dfrac{R_2[f,\mathbf{x}_0](\mathbf{x})}{\|\mathbf{x-x}_0\|^2}=0.
  \]
  Si, adem\'as, $f\in\mathcal{C}^3\implies\exists\:\mathbf{c}\in\Rn{n}$, en el segmento que une $\mathbf{x}$ con $\mathbf{x}_0$, tal que
  \[
    R_2[f,\mathbf{x}_0](\mathbf{x})=\sum_{i,j,k=1}^{n}\frac{1}{3!}\dfrac{\partial^3 f(\mathbf{c})}{\partial x_i \partial x_j \partial x_k}[\dots].
  \]
  \textcolor{red}{no se como escribir bien el resto}
\end{theorem}


\begin{definition} \textbf{Extremos locales y globales.} \label{def:extremos}
 \mbox{}
 
 Sean $f: A \subset \Rn{n} \to \R$ y $\mathbf{x}_0 \in A$. Decimos que:
 \begin{itemize}
  \item $f \text{ tiene un \emph{m\'aximo local} o \emph{relativo} en } \mathbf{x}_0 \text{ si } \exists\: \delta > 0 : f(\mathbf{x}_0) \ge f(\mathbf{x}),\;\forall \mathbf{x} \in B_{\delta}(\mathbf{x}_0) \cap A.$ 
  \item $f \text{ tiene un \emph{m\'inimo local} o \emph{relativo} en } \mathbf{x}_0 \text{ si } \exists\: \delta > 0 : f(\mathbf{x}_0) \le f(\mathbf{x}),\;\forall \mathbf{x} \in B_{\delta}(\mathbf{x}_0) \cap A.$
  \item $f \text{ tiene un \emph{m\'aximo global} o \emph{absoluto} en } \mathbf{x}_0 \text{ si } f(\mathbf{x}_0) \ge f(\mathbf{x}),\;\forall \mathbf{x} \in A.$ 
  \item $f \text{ tiene un \emph{m\'inimo global} o \emph{absoluto} en } \mathbf{x}_0 \text{ si } f(\mathbf{x}_0) \le f(\mathbf{x}),\;\forall \mathbf{x} \in A.$
 \end{itemize}
 Decimos que $f$ tiene un \emph{extremo} en $\mathbf{x}_0$ si $f$ tiene un m\'aximo o m\'inimo (absoluto o relativo) en $\mathbf{x}_0$.
\end{definition}

\begin{obs}
 Notemos que todo extremo absoluto es \emph{tambi\'en} un extremo relativo. La rec\'iproca de esta afirmaci\'on es falsa (un extremo relativo no tiene por qu\'e ser absoluto).
\end{obs}

\begin{theorem}\textbf{Condici\'on necesaria de extremo.} \label{teo:grad_nulo}
 Sea $f: A \subset \Rn{n} \to \R$ diferenciable en alg\'un entorno abierto de $\mathbf{x}_0 \in \interior{(A)}$. Supongamos que $f$ tiene un extremo en $\mathbf{x}_0$. Entonces
 \[
  \grad f(\mathbf{x}_0) = \mathbf{0}.
 \]
\end{theorem}

\begin{definition} \textbf{Punto cr\'itico.} \label{def:pto_crit}
 Sean $f: A \subset \Rn{n} \to \R$ y $\mathbf{x}_0 \in \interior{(A)}$. Decimos que $\mathbf{x}_0$ es un \emph{punto cr\'itico} de $f$ si $f$ no es diferenciable en $\mathbf{x}_0$ o $\grad f{(\mathbf{x}_0)} = \mathbf{0}$.
\end{definition}

\begin{definition} \textbf{Punto silla.} \label{def:pto_silla}
 Sean $f: A \subset \Rn{n} \to \R$ y $\mathbf{x}_0 \in \interior{(A)}$. Si $\mathbf{x}_0$ es un punto cr\'itico de $f$ y $f$ no tiene un extremo en $\mathbf{x}_0$, decimos que $f$ tiene un \emph{punto silla} en $\mathbf{x}_0$.
\end{definition}

\begin{theorem}\textbf{Criterio de la derivada segunda.} \label{teo:derivada_2da}

   Sean $A \subset \Rn{2}$ un conjunto abierto y $f: A \subset \Rn{2} \to \R$ una funci\'on de clase $\mathcal{C}^2$. Sea $(x_0,y_0) \in A$ y supongamos que $\grad f(x_0,y_0) = \mathbf{0}$. Entonces,  
   \begin{enumerate} %[I.]
      \item si $\det(\mathbf{H}_f (x_0,y_0)) > 0$ y
      \begin{enumerate} %[(a)]
          \item $\pdv[2]{f}{x} (x_0,y_0) > 0$, $f$ tiene un m\'inimo local en $(x_0,y_0)$;
          \item $\pdv[2]{f}{x} (x_0,y_0) < 0$, $f$ tiene un m\'aximo local en $(x_0,y_0)$;
      \end{enumerate}
      \item si $\det(\mathbf{H}_f (x_0,y_0)) < 0$, $f$ tiene un \emph{punto silla} en $(x_0,y_0)$.
   \end{enumerate}  
\end{theorem}

\begin{obs}
  \textcolor{red}{decir pq se puede evaluar tanto $f_{xx}$ como $f_{yy}$}
\end{obs}

\begin{theorem}\textbf{de los valores extremos de Weierstra\ss.} \label{teo:weier}
  Sean $A \subset \Rn{n}$ un conjunto cerrado y acotado y $f: A \subset \Rn{n} \to \R$ una funci\'on continua. Entonces $f$ alcanza un m\'aximo y un m\'inimo absoluto en $A$.
\end{theorem}

\begin{theorem} \textbf{M\'etodo de multiplicadores de Lagrange.} \label{teo:lagrange}
    Sean $f:U \subset \Rn{n} \to \R$ y $g:U \subset \Rn{n} \to \R$ funciones de clase $\mathcal{C}^1$. Sea $\mathbf{x}_0 \in U$ y sea $c = g(\mathbf{x}_0)$. Sea $S$ el conjunto de nivel $c$ de $g$. Supongamos que $\grad g (\mathbf{x}_0) \ne \mathbf{0}$. Si $f|_S$ (que denota a la restricci\'on de $f$ a $S$) tiene un extremo local en $\mathbf{x}_0$, entonces $\exists\: \lambda \in \R$ tal que
    \[
     \grad f (\mathbf{x}_0) = \lambda \grad g (\mathbf{x}_0).
    \]
\end{theorem}
